% META: {"id": "1NSI-RES-EX-008", "chapitre": "1NSI-RESEAUX", "type_objet": "exercice", "capacites": ["1NSI-RESEAUX-C2"], "parcours": 2, "competences": ["analyser", "modeliser"], "duree_min": 8, "mode_creation": "ex_nihilo", "sources_inspiration": [], "fichier_tex": "chapitres/1NSI-RESEAUX/exercices/1NSI-RES-EX-008.tex", "corrige_tex": "chapitres/1NSI-RESEAUX/corriges/1NSI-RES-CO-008.tex", "status": "approved", "capacites_codes": ["C2"], "methodes": ["M1"]}
% BEGIN-VERIFY
% reseau = {
%     "Salle1": ["Switch1"],
%     "Salle2": ["Switch1"],
%     "Switch1": ["Salle1", "Salle2", "Switch2"],
%     "Switch2": ["Switch1", "Salle3"],
%     "Salle3": ["Switch2"],
% }
% def peuvent_communiquer(reseau, depart, arrivee, visites=None):
%     if visites is None:
%         visites = set()
%     if depart == arrivee:
%         return True
%     visites.add(depart)
%     for voisin in reseau.get(depart, []):
%         if voisin not in visites:
%             if peuvent_communiquer(reseau, voisin, arrivee, visites):
%                 return True
%     return False
% assert peuvent_communiquer(reseau, "Salle1", "Salle3") is True
% reseau2 = dict(reseau)
% reseau2["SalleIsolee"] = []
% assert peuvent_communiquer(reseau2, "Salle1", "SalleIsolee") is False
% END-VERIFY
\begin{exercice}{1NSI-RES-EX-003}{2}{12}
Le réseau d'un établissement scolaire relie trois salles à deux commutateurs :
\begin{python}
reseau = {
    "Salle1": ["Switch1"],
    "Salle2": ["Switch1"],
    "Switch1": ["Salle1", "Salle2", "Switch2"],
    "Switch2": ["Switch1", "Salle3"],
    "Salle3": ["Switch2"],
}
\end{python}
\begin{enumerate}
  \item En utilisant \lstinline{peuvent_communiquer} du cours, vérifier que \lstinline{"Salle1"} et \lstinline{"Salle3"} peuvent communiquer, bien qu'elles ne soient reliées à aucun commutateur commun directement.
  \item Ajouter une machine \lstinline{"SalleIsolee"} qui n'est reliée à \textbf{aucune} autre machine du réseau (dictionnaire avec une liste vide). Vérifie qu'elle ne peut communiquer avec aucune autre salle.
  \item Pourquoi la fonction marque-t-elle les sommets \og visités \fg{} pendant sa recherche ? Que risquerait-il de se passer si ce réseau contenait un cycle (par exemple si \lstinline{Switch1} et \lstinline{Switch2} étaient reliés par \textbf{deux} câbles distincts, modélisés comme une arête supplémentaire) et que cette précaution était absente ?
\end{enumerate}
\end{exercice}
